% pass-fix-f203-alg13-device-keygen-k4 — FIPS 203 Alg.13 ML-KEM-768 KeyGen 验收探针 (k=4)
% 编译：bash ../../scripts/xelatex-clean.sh pass-fix-f203-alg13-device-keygen-k4-实现方案-customspec.tex
\documentclass[UTF8,a4paper,11pt]{ctexart}
\usepackage{amsmath,amssymb}
\usepackage{booktabs}
\usepackage{geometry}
\usepackage{hyperref}
\usepackage{longtable}
\usepackage{array}
\usepackage{seqsplit}
\geometry{margin=2.2cm}
\newcolumntype{L}[1]{>{\raggedright\arraybackslash}p{#1}}
\newcommand{\apiname}[1]{\texttt{\seqsplit{#1}}}
\newcommand{\dirname}[1]{\texttt{\seqsplit{#1}}}
\hypersetup{colorlinks=true,linkcolor=blue,urlcolor=blue}

\title{pass-fix-f203-alg13-device-keygen-k4\\FIPS 203 Alg.13 ML-KEM-768 PKE KeyGen 实现方案（k=4）}
\author{ascendc-tests 验收探针（pass- 终态）}
\date{2026-06-28}

\begin{document}
\maketitle

\begin{abstract}
本文档描述 \textbf{验收探针} \dirname{ascendc-tests/pass-fix-f203-alg13-device-keygen-k4/}：
在 \textbf{Atlas A2 / Ascend910B4} 上以 AscendC 实现 FIPS 203 \textbf{Algorithm 13}（ML-KEM-768，$K=4$）\textbf{PKE KeyGen} 全设备路径。
\textbf{生产 I/O}：\texttt{input/} 仅 \texttt{SEED\_D} + NTT LUT；\texttt{output/} 仅 \texttt{ek\_PKE}(1568B)/\texttt{dk\_PKE}(1536B)。
\textbf{Host 入口}：单进程 \texttt{ascendc\_keygen\_bbit}（\texttt{cmake/keygen}）；\textbf{2 次 device launch}（prep $\vert$ compute+行21融合）。
源码 \textbf{vendored} 于本目录 \texttt{prep/}、\texttt{compute/}；仅共享 \texttt{library/shared} 的 SHAKE/Keccak。
交付示例对照：\dirname{examples/incubating/exp-mlkem-f203-pke-keygen-k4/}（用户 I/O 契约相同，本探针为验证基线）。
全目录源文件含 \texttt{@probe} 注释块（见 \dirname{scripts/inject\_probe\_code\_comments.py}）。
\end{abstract}

\tableofcontents
\newpage

\section{数学语义（Alg.13，$K=4$）}

\subsection{输入输出}
\begin{itemize}
  \item 输入 \texttt{SEED\_D}：\texttt{uint32}；探针约定 \texttt{DerandFromSeedD}（与 liboqs KAT 对齐）。
  \item 输出 \texttt{ek\_PKE}：1568 B；\texttt{dk\_PKE}：1536 B。
  \item 中间 GM：单次 \texttt{ascendc\_keygen\_bbit} 进程内驻留；\textbf{禁止}默认写入 \texttt{input/}/\texttt{output/}。
\end{itemize}

\subsection{流水线}
\begin{enumerate}
  \item $G(d\Vert k)\to(\rho,\sigma)$（prep 设备 Keccak）。
  \item 行 3--7：16$\times$ SampleNTT$(\rho)\to\hat{A}$。
  \item 行 8--15：PRF$(\sigma)$+ CBD$_{\eta_2}\to\hat{s}$（\texttt{src}）。
  \item 行 16--20：2s1e NTT + MMAD + 行18内积 + ByteEncode$_{12}$。
  \item 行 21：\texttt{ek\_PKE = ek\_polyvec $\Vert$ $\rho$}（\texttt{F203\_KEYGEN\_EK\_PKE=1}，与行16--20同次 Launch 2）。
\end{enumerate}

\section{生产编排与 I/O}
\label{sec:io}

\begin{longtable}{@{}L{3cm}L{10cm}@{}}
\toprule
步骤 & 行为 \\
\midrule
准备 & \texttt{scripts/prepare\_production\_input.py} $\to$ \texttt{input/seed\_d.bin} + LUT \\
构建 & \texttt{cmake/keygen} $\to$ \texttt{ascendc\_keygen\_bbit} \\
运行 & 2 launch；GM 指针传递；仅落盘 \texttt{ek\_pke.bin}、\texttt{dk\_pke.bin} \\
清理 & \texttt{\_keygen\_scrub\_output} 删除其它 \texttt{output/} 文件 \\
SIM & build \textbf{后} \texttt{sim\_env\_export}（WSL stub）；\texttt{camodel\_sim\_log.sh} \\
\bottomrule
\end{longtable}

\textbf{禁止}作为生产接口：\texttt{compute\_io/} 磁盘 staging、\texttt{run\_prep}+\texttt{run\_compute} 分段默认路径、向 \texttt{input/} 写 \texttt{a\_hat/src/rho}。

\section{Launch 与 AI Core}
\label{sec:aicore}

\begin{longtable}{@{}L{2.2cm}L{3.5cm}L{2cm}L{5.3cm}@{}}
\toprule
Launch & 内核 & blockDim & SIM profile（910B4） \\
\midrule
1 & \texttt{f203\_keygen\_prep} & 2 & \texttt{type=AIV}；0 AIC cycle；2 AIV block（block0 $\gg$ block1） \\
2 & \texttt{mmad\_custom} & 1 & \texttt{MIXAIC}；1 AIC + 2 AIV；\texttt{core\_list=[0]} \\
\bottomrule
\end{longtable}

\begin{itemize}
  \item 两次 launch \textbf{串行}占用 \textbf{1 颗 AI Core}（非双核并行）。
  \item prep 当前实现：\texttt{F203\_AHAT16\_BLOCK\_DIM=2} 但 Â 由 block0 \textbf{串行}两片（SIM 曾 block1 并行写 GM maxdiff$\approx$3300）；block1 主要在 P-04 同步。
  \item CPU \texttt{[SUCCESS][AIC\_x]} 为 tikicpu 芯片拓扑 artifact；\textbf{以} \texttt{sim\_log/profile\_*.toml} \textbf{为准}。
\end{itemize}

\section{目录与构建}

\begin{verbatim}
pass-fix-f203-alg13-device-keygen-k4/
  main_keygen.cpp          # 生产 host（默认 run.sh）
  cmake/keygen/           # 统一 prep+compute 内核库
  prep/{ahat,alg7,presample,alg8}/   # vendored 行 3-15
  compute/                # vendored vec-k4-v2 行 16-21
  scripts/prepare_production_input.py
  run.sh
\end{verbatim}

遗留分段：\texttt{cmake/prep}、\texttt{cmake/compute}、\texttt{main\_keygen\_prep.cpp}、\texttt{main\_compute.cpp}（调试专用）。

\section{数据与 Pipe 同步（prep）}

详表：\dirname{PIPE\_SYNC\_EVAL.md}；原理：\dirname{docs/notes/F203-KeyGen-prep-Pipe细同步技术总结.md}。

prep 与 compute 之间 \textbf{无设备 barrier}；靠 Launch 1 结束 + Host Launch 2 顺序。

\section{验证}

\begin{verbatim}
cd ascendc-tests/pass-fix-f203-alg13-device-keygen-k4
bash run.sh -r cpu -v Ascend910B4
bash run.sh -r sim -v Ascend910B4
KEYGEN_VERIFY=1 bash run.sh -r cpu -v Ascend910B4
bash kat_liboqs_vs_ascendc.sh
\end{verbatim}

\section{性能（SIM 910B4，生产单入口，2026-06-28）}

Total tick \textbf{886801}（task0 prep $\approx$806k + task1 mmad $\approx$81k）。

\end{document}
